\chapter{生涯规划}
\SourceParagraph{南赫学院至今还没有第一届本科毕业生，如果要我对大家的就业去向做个预测，我相信会是“泻水置平地，各自东西南北流”。或许现在并不是写这个章节的好时候，毕竟我们中还没有人在现实中经历过就业升学。不过抛开可能潜在的功利因素，我知道有许多同学正在因为未来就业问题而迷茫徘徊、上下求索。这一章节我们注定做不到面面俱到，错谬之处在所难免，只希望能分享一些信息差，尽我们的绵薄之力来抛砖引玉。如果后面的内容使你受益，也期望你加入到我们的项目中，分享自己的感悟。}
\SourceParagraph{不论你现在的偏好如何，大二下的暑假最好能确定自己未来的方向。这里说的方向是指本科毕业之后将要直接就业、保研、考研抑或是出国，是否适合并考虑选择学术界，是否将会投身工业界种种。原因如下：首先保研看前五个学期的学分绩，这个时间点已经基本能够确定自己能否保研了，如果不能保研的话，从大二下的暑假开始还有一年多的时间可以准备考研，准备时间比较充足。想要做学术从暑假可以开始寻找合适的导师做长时间的科研，想要进工业界可以开始准备实习，想要出国现在开始准备刚好。不论你自己选择哪个方向努力，我都有三点建议：1，你需要和家庭协商。关于未来发展道路的任何决定都最好能得到家庭的支持。2，你最好有一个底线思维。倘若考研考不上怎么办，家庭会支持你二战吗/转专业失败怎么办，最差自己会干什么，自己又能否接受这个最差的结果。3，不要闭门造车。想出国就和已经出国的人聊，想跨专业就努力融入其他专业的环境，收集信息很重要，善于寻求帮助也是一个难能可贵的品质。}
\SourceParagraph{每个人性格不同，想法不同，选择也不同，一件重要的事是选择本身没有高下之分。上述提及的路不管那一条，都已经有大气或难喝的学长学姐涉猎。有了下面的内容，相信你们会比我们和他们走得更远。}
\section{境内升学}
\subsection{大气保研}
\SourceParagraph{事先声明，目前南赫学院第一届本科生还没保研，所以一切都只能参考大气院的经验，以及学院领导/老师的说辞。并且 2025 年开始，过去的夏令营制度出现了重大变化，目前尚不清楚这会对保研，尤其是外保产生何种影响。}
\SourceParagraph{保研，也就是推荐免试研究生，是常用升学手段之一。通常根据专业和学校进行排列组合，可以分为本校本专业，外校本专业，本校跨专业，外校跨专业。保研有两个关键步骤，一个是获得本院的推免资格，一个是目标学校接收。下面围绕这两个流程来进行阐述。}
\SourceParagraph{推免资格是分配到学院的，通常是按一定人数比例分配，虽然名额每年会有浮动，但是往年情况也可以作为参考。大气院通常是 50\%左右，但是他们是有拔尖计划的，所以不好一概而论，而南赫学院目前没有分班（23 级），也没有对学生作政策上的区分，所以内部评价标准是均质的。}
\SourceParagraph{据目前得到的消息，}
\SourceParagraph{南赫学院的保研名额是学生总数的 40\%左右，尚不清楚该比例会如何发展。}
\SourceParagraph{如何获得推免资格？在大三下的暑假（八月底九月初）学院会发布正式的方案，通常是按照推免学分绩和学生自主申请的情况分配推免资格。这里的推免学分绩主要是考察学位课（准出课），但是也不完全等于学位课，哪些课算哪些课不算，该怎么算，取决于学院。南赫学院的必修课程比较抽象，目前还不能得知学院会如何考虑不同模块的必修课。课程成绩按照第一次参加考试的得分记，所以千万不要想着补考刷分对保研有用！如果是缓考的话，只会计算 90\%的分，但是如果是因公缓考则不必，具体的方案还得等学院推出。除此之外，还存在极少数的加分，如发表论文，各种竞赛/活动获奖（如南大基础学科论坛，挑战杯，“互联网+”，数学建模国赛/美赛等），细则可能会有变化，还有一些审核流程和具体要求，并且每个学院要求可能不一致，仅供参考，仅供参考，仅供参考！！！}
\begin{figure}[H]
  \centering
  \includegraphics[width=0.92\textwidth,height=0.78\textheight,keepaspectratio]{figures/p25\_06.png}
\end{figure}
\begin{figure}[H]
  \centering
  \includegraphics[width=0.92\textwidth,height=0.78\textheight,keepaspectratio]{figures/p25\_07.png}
\end{figure}
\SourceBlankLine[2]
\SourceParagraph{按照上述流程得到最终推免学分绩，然后按照提交了推免申请的人进行排序。所以说，想要保研，最主要的是保证学位学分绩排名，靠比赛能加的分非常有限，远达不到扭转局势的程度，不如把握住能把握的分，尤其是学分多的课，比如大一、大二上的数学、物理、英语课程。}
\SourceParagraph{获得了推免资格，或者预计自己能够获得推免资格，接下来就要考虑去向了。}
\SourceParagraph{保本校本专业是最轻松的。联系老师宜早不宜迟，因为大部分老师选择学生并不是优胜劣汰，而是先来后到，提前锁定老师，老师答应的人够名额之后，一般也就不会答应别人了。从这个意义上，保本校本专业不需要养鱼，找到愿意收你的老师就可以了，养鱼反而日后尴尬。可能比较有挑战的是一些大牌的老师，比如院士组，他们会对学生有一些要求，但是也不全是绩点导向的，打铁还需自身硬。大部分老师都有至少一位硕士名额的，但是有些老师不喜欢硕士所以不招，这点要问清楚。}
\SourceParagraph{如果要保外校，就没有这么轻松了。按往年来说，大三下的暑假开始，各个单位就会陆续开展夏令营，如果想外保一定要争取参加夏令营，并且争取在夏令营中好好表现，拿到优营。对某些学校，比如北大来说，优营就是铁 offer。但是没有拿到优营也有机会，因为总有放弃的人。虽然流程上没有参加夏令营也可以申请，但是夏令营作为线下考核、和导师交流、展示自我的绝佳平台，大多数情况下还是参加了才比较有希望。理论上参加夏令营没有学分绩的硬性要求，但是肯定越高越好，每年情况也不一样。以北大为例，}
\SourceParagraph{21 级及之前的，报名了基本都能入营，但是 22 级有很多人报，但是情况不太理想。这时候就不得不提 2025 年开始的新政策了。大多数学校不开展夏令营了，疑似是收到了上面什么通知，但是还是以开放日的形式开展了类似的活动。事态会如何发展，我们持续关注。保外校可以不急着联系老师，因为夏令营可以再联系，但是也可以先联系着，利用各种会议、暑校的机会去接触老师。}
\SourceParagraph{本校可以不参加夏令营直接参加预推免的考核，通常是面试，外校的在夏令营/开放日中会有另外的考核。And vice versa，考虑保外的时候我们也一样，需要夏令营的考核。不过由于预推免名单确定，面试基本就是走个过场，只要别太拉跨就行。}
\SourceParagraph{推免材料里会有专家推荐信这一项，通常就是找自己的授课老师或者指导老师就行，不算关键。}
\SourceParagraph{关于如何选择导师，首先肯定是老师的研究方向和自己匹不匹配，这点可以通过查看老师简介和论文、跟老师发邮箱甚至线下约时间聊一聊了解，可以结合自己未来的规划（科研？就业？）寻找适合自己的老师。其次要看一下老师的帽子和年龄，年纪比较大的老师可能不会很有空辅导你，比如谈院士的学生一年就见四次谈校，往往实际还是让师兄师姐或者小老板带，他们能够提供的指导少，但是往往有比较好的 idea，能够给出比较好的研究方向上的建议，具体研究的话很多时候还得是自力更生。年轻老师就能给比较多的指导，相应地也会比较 push，如何选择还是要看自己。可能在读研读博的过程中，老师的帽子就变多了，职称也变高了，也变忙了，要用发展的眼光看待。最后还要关注一下老师的性格什么的，毕竟是要跟 3-5 年（南赫硕的话就是 2 年）的，好不好相处啥的还是挺重要的，这些可以去问老师的学生。南大的老师还是很好相处的，外校的话，可以找找机会接触，利用刚才说的那些，会议、暑校，甚至是讲座之类的。一些更具体的信息可以用邮箱联系一下他们的学生，或者发帖（表白墙、知乎、小红书等等）捞。}
\SourceBlankLine[2]
\subsection{大气考研}
\SourceParagraph{约稿中\textasciitilde{}\textasciitilde{}}
\newpage
\section{境外升学}
\subsection{留学1——留学基础情况介绍QwQ}
\SourceParagraph{这是一篇根据作者自己搜索的资料和自身目前的经历总结的留学经验分享。虽然本人一直在自行规划留学，收集信息和超过10+中介斗智斗勇了，且采访了多位学长/学姐/教授。目前本人准备申请北美对口专业的研究生，但是受限于目前才大二，也仅申请大气对口专业，在这里仅做一个抛砖引玉的作用。这篇算是南赫特供版demo。}
\SourceParagraph{不要盲目选择出国，特别是读硕士，真的很贵。}
\subsubsection{一、出国留学的实用性和目的}
\SourceParagraph{这一段主要是让大家对留学祛魅。让我们扒光留学华丽的袍子，看清一地的鸡毛和爬满的虱子。}
\SourceLabel{现在出国留学最适合的是这几类人：}
\SourceParagraph{第一类，资源支持型。}
\SourceParagraph{家里可以轻松承担留学全周期费用（但是博士生是有funding的，之后会讲），且具备职业安排能力。简而言之就是各类公子小姐。这类人相比于知识，更需要的海归的人设和圈子，或者说就是去体验生活。}
\SourceParagraph{第二类，学历镀金型。}
\SourceParagraph{对于那些第一学历不太好，需要硕士来提升学历的，或者保研考研考公都失败没地方可去的，读一个硕士还是很有必要的。因为有一些企业是卡学历必须要是硕士，而且必须要QS前50/100。同时有一个QS高学历方便落户（注：请参考不同地方和公司的政策）。这一类留学的时候请一定要安排好自己的实习和找工作，留学是作为找到好工作的手段，不是目的。}
\SourceParagraph{第三类，海外发展型。}
\SourceParagraph{想要润国外现在最便宜的方式就是留学，分为找工作和找教职。留美国仍然主流，但政策逐渐收紧。如果你真的不喜欢国内的氛围，能早出去早出去，为此读一个硕士当跳板都不要紧。这边通过举例说明境外就业方案：}
\SourceParagraph{转码 -> 纽约大学CS硕士（要一定的前期话费在学费生活费上）+当地实习-> OPT工签 -> H1B抽签 —> 拿到绿卡一直当码农}
\SourceParagraph{博士获全额资助（Cover学费+生活费）→ 留校研究/工业界就业}
\SourceParagraph{（注：这条路其实也很难走，有可行性但是要看命和概率。可以自己去知乎搜集资料）}
\SourceParagraph{当然你也可以不在国外发展，可以利用留学生身份拿到3-5年的工签，在当地的公司实习一段时间再回国发展。但是走这条路的同学一定要注意自己读的硕士项目，不要读太水的。这条路还有一个好处是可以转专业（大气可以转统计/DS/转码）}
\SourceParagraph{第四类，学术型（这一类可能南赫多一点）}
\SourceParagraph{立志这辈子做科研的同学可以申请博士（或者先硕士再博士）。这样评教职会更有竞争力，且眼界开阔，职业选择开阔。可以在国外或者国内找教职继续做科研，也可以找工作去工业界。虽然说现在学历贬值，但是一个海外学校博士含金量依然足够，不太会遇到硕士找工作的困境（特别是现在的ai大环境下，强理工科博士吃香，短时间内难以被ai取代）。读博是有funding的，是可以cover生活开销的。但是申请国外直博是非常难的，现在读国内博士也不一定比出国差，国内有一些领域科研水品是比国外高的。（作者向杜灵杰教授和傅豪教授求证过，如果没有强大的个人愿景或者申请到了超级牛校的博士，综合考虑出国读博和清北读博各有优劣。就算想去国外读博也可以现在国内读硕士当跳板，申请结果更好）而且准备海外申直博注意点很多（后两篇会分散讲到，特别是第二篇基本上就可以认为是博士申请材料指导，因为我是按照这个条件准备的，虽然最后大概率是去读硕士，但是这个条件最高，向下兼容，且部分硕士项目例如普林Mfin，斯坦佛ICME录取要求和难度不比博士简单）。}
\SourceLabel{留学的目的和收获：}
\SourceParagraph{最本质的一点就是经历和眼界。国外的氛围，异国的文化，当地的氛围，遇到的人，见到的事，接触的信息，拥有的机会，在准备和留学过程收获的能力。这个是最玄妙的一点，因为它既无法用金钱衡量，也很看个人操作（有的留学老鼠人就是天天公寓，学校，超市三点一线。完全没有利用好留学平台）}
\SourceParagraph{再其次就是受到的教育，关于国外的教学水品，知乎上有很多回答，不再赘述。我的理解是中国的大学收到了90分的学生，经过教育后后学生成长到了95分，而国外大学的学生入学时期可能只有80分的水平，但是经过教育后也可以到95分。（可以看《上海交通大学学生生存手册》对国内高等教育的评价，简单来说就是已经奔溃。作者本人这次去清北访问了一下，更加感到国内教育体系的畸形还有国内内卷压抑的环境）当然学校只是提供氛围环境，对于真正天才和狠人，所有教育都黯然失色。}
\SourceParagraph{最后就是connection和资源，能去国外读书的人都是经过筛选的，圈层是更好的；而且虽然现在主体是东升西落，但是国外还是有很多的工业界和学术界的资源，这是一个圈层和派系（学术界是绝对有派系的）的问题（这次作者去北京参加学术会议感受到，那些教授都是互相照应的、有圈层的，不是同一个学校/派系的人很难与他们有联系）。且很多教授真的很厉害，自己拥有公司或者认识很多公司老板。（注：请注意，所有资源和人脉都是在你自身也有价值的情况下才能调动起来，请勿盲目相信人脉资源，和一个喜欢吹牛皮的中年老登一样。本质上是资源交换和互利共惠）}
\SourceParagraph{最后一点才是一张学历。现在学历真的是在贬值（博士及以上学历贬值幅度低），大家都只是卡一个985/211或者qs50/100。真正以后重要的是一个人的能力，找工作时也是面试和笔试更加重要，学历只是一个筛选的敲门砖。而且对于大多数情况和大部分安安稳稳找工作度日的老百姓，南大本硕已经很够用了，再高的学历边界效应过于突出。（少部分情况就是如果你想去顶级量化公司，顶级ai公司。清北，姚班，hypsm，牛剑都是你的竞争对手，读一个好一点的硕士博士都只能弥补一部分）}
\SourceBlankLine[1]
\SourceParagraph{在这个学历贬值，留学难以金钱上回本的时代，请大家结合自身情况，看清楚自己究竟想要什么再选择是否留学——你能否且应该为学历烧掉一套房？而这一份学历究竟能带来什么？（你也不希望花几年青春和几百万人人民币换来一张废纸吧）}
\subsubsection{二、申请的大致介绍}
\SourceParagraph{国外大学是申请制度的，大家要习惯没有分数线和确定性，就算你材料完美也有可能被拒，真正决定拍板的再也不是投档线和分流排名，而是招生组和教授（所以有一些人为因素，包括connection和推荐信显得很重要）。需要的参考材料有以下这几个：学校背景，GPA，语言成绩，科研/竞赛/实习经历，论文成果，推荐信，文书。}
\begin{center}
\small
\setlength{\tabcolsep}{3pt}
\renewcommand{\arraystretch}{1.55}
\begin{tabular}{|>{\Centering\arraybackslash}m{0.14\textwidth}|>{\Centering\arraybackslash}m{0.30\textwidth}|>{\Centering\arraybackslash}m{0.18\textwidth}|>{\Centering\arraybackslash}m{0.27\textwidth}|}
\hline
\multicolumn{4}{|c|}{留学要求逻辑总结图} \\
\hline
  & 院校背景 & & \\
\cline{2-2}
  & 专业对口度 & 学术背景 & \\
\cline{2-2}
录取原则 & 平均分 GPA & & \\
\cline{2-4}
  & 软背景（实习/科研/竞赛等） & & \\
\cline{2-2}
  & 标化考试（雅思/GRE/Gmat） & 非学术背景 & \shortstack{有具体内容的实习\\有收录结果的科研\\有高度匹配的经验} \\
\cline{2-2}
  & 文书质量/推荐信质量 & & \\
\hline
\end{tabular}
\end{center}
\SourceBlankLine[2]
\SourceParagraph{这些材料是什么，有什么作用，要怎么准备会在下一篇详细论述。}
\SourceLabel{申请的大致流程：}
\begin{figure}[H]
  \centering
  \includegraphics[width=0.92\textwidth,height=0.78\textheight,keepaspectratio]{figures/wq\_image2.png}
\end{figure}
\SourceParagraph{下一篇会有更加详细和适合南赫宝宝体质的时间规划。}
\SourceBlankLine[1]
\subsubsection{三、留学的细分}
\SourceParagraph{注：这一段有的\RareLuo{}嗦，而且这一段很多的话对于一个刚准备留学的人来说过于概括和抽象，建议了解一部分信息后再来观看。坚定留学的一定要把后面两篇一起看了才能有体会。且强烈建议自己再去多找资料交叉验证，结合具体数据体验——比如录取率，录取bg，毕业年薪，学费开销。}
\SourceParagraph{留学的选择具体在地区上可分为：英国，香港，新加坡，澳洲，新西兰，美国，加拿大，日本，韩国，马来西亚等。这边从申请体制出发，去掉一些不太热门的留学地点，分为英体制（英国，香港，新加坡）和北美体制（主要是美国，加拿大还是有点冷门了）介绍。澳洲的低门槛导致学生圈子和水平略显抽象，现在被打上了水校的标签，而且现在成本逐年上升，不做介绍（澳洲和英国【牛剑除外】有点畸形，硕士门槛和含金量低，但是博士有大牛教授）；马来西亚大学主要受众不是985学生；日本韩国过于小众而且申请制度作者没有了解OvO。在申请的学位方面，这边介绍硕士和直博(不着重介绍小众的二硕，和部分项目特有的硕士转博士)}
\SourceParagraph{英国体制的学校在硕士申请时更加在申请时更加看重硬背景，即本科院校+GPA+语言成绩就基本上决定了你可以申请到的学校。去英国留学的主要动力是：}
\SourceListItem{1 英国体制的学校一般来说就读时间短，时间成本低。}
\SourceListItem{2 香港和新加坡华人社会较好融入，且离家较近。}
\SourceListItem{3 总体费用相对于美国便宜。}
\SourceListItem{4 英国体制的学校qs都偏高，且申请难度显著比美国低（GPA占大头）。}
\SourceParagraph{北美的学校则更加折腾一点，硕士申请时软硬背景基本五五开，要申请到好学校要当五边形战士。（会在下一个篇章具体讲解申请的材料，到时候是根据美国的申请事项来讲的，因为这是向下兼容的，要是你连五边形战士都能当好，别的学校的申请不在话下）。去北美留学的主要动力是：}
\SourceListItem{1 含金量和认可度高。985+美硕，如果是回国发展基本是学历顶配。但是体制内的企业可能更加认可国内研究生。}
\SourceListItem{2 不是笔者崇洋媚外，但是目前来说美国还是有一些值得我们学习和借鉴的方面的。希望同学们学完成归来，报效祖国，建设共产主义社会。}
\SourceListItem{3 可以留在当地，当润人。前面有提到过具体的操作流程。但在这我不做过多宣传，好处坏处都有，大家自行定夺。值得一提的是要是选择留美，前期留学的话费是可以在工作内5年挣回来的，部分高薪专业真的可以达到年薪百万人民币；博士就更加不用说了，基本是免费留学，甚至可以电表倒转赚钱。（但是考虑物价换算，人民币和美元的换算绝对不能仅按照7:1。这边给出一个参考，生活消费大致3:1，电子产品/奢侈品/汽车9:1，房价另说）}
\SourceListItem{4 美国的学校比英体制学校好转专业，去年大气科学毕业生中就有两个去NYU学计算机的，如果大家不喜欢大气科学和其对口专业，留美换赛道是十分可行的，但是要接受申请的学校档次会低一档，大学本科也要同时学好两门专业，更加刻苦。}
\SourceBlankLine[1]
\SourceParagraph{硕士在国外是要自掏腰包就读的（很难拿到奖学金）。一些水学校的水项目来近年来有变成付费揽钱的趋势，所以也有人调侃为水硕，而且含金量越来越低，仅读硕士对于一般人来说意义不大。但是有含金量的学校的硕士项目是难申请的（牛剑和美国大藤），硕士申请一般是招生办审核，招生时筛选的标准大致是有学习能力的人。就读分为授课型和研究型，可以简单粗暴的理解为授课型是工作导向，研究型是科研导向。一般选择读研究型的是为了曲线救国，以后再申请博士或者直接转成博士（注：有的好学校的研究型硕士项目申请的要求和难度是向博士申请看起的）。具体请按照各个项目学校官网具体的培养计划，这边仅是让大家有初步认知。硕士不同的项目之间bar差很多（bar可以理解为录取难度和含金量），要结合自己的情况来抉择。一般来说没有确切的规划都可以先考虑的申请硕士，边读硕士边探索边考虑未来。}
\SourceParagraph{直博在国外一般是有funding既教授发的工资的，是属于教授自掏腰包培养科研人才。所以申请难度极大，含金量上更加看重你跟随的教授和学科排名。现在国外普遍削减科研经费，很多博士项目减少招生甚至不招生，近几年申请难度陡增。我们说的offer狭义上就是指拿到博士录取，因为offer是不用自掏腰包的。硕士录取我们一般叫AD，仅代表你被录取了，当然硕士也有可能拿到全额奖学金的offer。博士招生一般来说来说会有教授组的亲自面试，而且要提前找中意契合的教授套辞联系。博士招生看的是你和教授课题组研究方向的适配度以及你的科研能力，毕竟教授是来培养科研人才，等着你来发paper并且在学术界深耕的，材料上更加注重科研经历和科研产出。如果你能上一个好学校的博士且自身有能力，是少数现在可以提升阶级的方法，教职/大厂/华尔街，真的是年薪百万。如果博士申请难度太大，可以先考虑在国内/外先读研究生，再申请博士曲线救国，一般来说这样能申请到的学校能提高一个档次，但是这样会有时间成本（直博基本要读5年+，先读硕士再读博士只会更久），但一般申请结果会更好看。读博士这个选择一定要慎重，首先时间成本高，而且读博士很 痛 苦，关于这点可以去看《THE PH.D. GRIND》，一位斯坦博计算机直博博士生美国人写的，我们一般称之为“PH.D.之殇”。但是只要有意志力任何人都是可以读博的，因为只要你能申请上博士就代表你能力没问题，剩下的交给汗水和时间。值得一提的是有的学校的有些科目是没有硕士只能申请博士的（比如普林斯顿的大气科学只有博士点），有的学校的项目硕士是可以转成博士的（比如斯坦佛的Civil and Environmental Engineering）。如果你读了博士但是想中途辍学，大部分学校可以在你达成要求的情况下给你“关怀补贴”——读满两年后退学给予硕士学位。}
\SourceParagraph{当然留学中不同的学科也是不同的，文科重视实践活动，理科注重竞赛科研，商科注重实习经验，医学生一般不留学（因为不如协和4+4（bushi）。其实是因为国内和国外的医学教育体系不一致，顶多当访问学者学习交流）。这边都是依照一个理工科视角展开的。}
\SourceParagraph{以下是这一part的总结}
\begin{center}
\small
\setlength{\tabcolsep}{2.5pt}
\renewcommand{\arraystretch}{1.22}
\begin{tabular}{|>{\Centering\arraybackslash}m{0.20\textwidth}|>{\Centering\arraybackslash}m{0.24\textwidth}|>{\Centering\arraybackslash}m{0.45\textwidth}|}
\hline
体系 & 代表地区 & 申请侧重点 \\
\hline
英联邦体制 & 英国/香港/新加坡 & 硬背景导向（院校背景 + GPA + 语言） \\
\hline
北美体制 & 美国/加拿大 & 综合素质导向（科研/实习/文书 软实力支撑） \\
\hline
其他选项 &  &  \\
\hline
\end{tabular}
\end{center}
\begin{center}
\small
\setlength{\tabcolsep}{2.5pt}
\renewcommand{\arraystretch}{1.22}
\begin{tabular}{|>{\Centering\arraybackslash}m{0.17\textwidth}|>{\Centering\arraybackslash}m{0.27\textwidth}|>{\Centering\arraybackslash}m{0.18\textwidth}|>{\Centering\arraybackslash}m{0.27\textwidth}|}
\hline
类型 & 授课型硕士 & 研究型硕士 & 直博 \\
\hline
目标 & 就业导向 & 学术过渡 & 学术深造 \\
\hline
费用 & 自费较多 & 部分资助 & 全奖 \\
\hline
招生 & 委员会 & 教授参与 & 教授主导 \\
\hline
核心 &  & 研究潜力 & 课题匹配度 \\
\hline
\end{tabular}
\end{center}
\SourceBlankLine[1]
\SourceParagraph{这边仅仅是作为一个引子带大家了解留学，后面两篇会有详细的介绍和指导（第二篇是申请的详细介绍，第三篇是有关留学中介）}
\SourceBlankLine[1]
\subsection{留学2——申请材料QwQ}
\SourceParagraph{选择看到这里的同学，我可以认为你是坚定的准备留学了，所以我会把我知道的基本毫无保留的告诉你（除了一些敏感数据和敏感资料）。如果你暂时还没有定夺或觉得信息量太大，可以大一下再来看，大一的寒假里选择一个交换项目体验一下国外学校的氛围，且自己再收集一些资料辅助判断，但请务必稳住你的GPA（作者就是因为大一上没有准备留学，学习摆烂。微积分只有89分，大化更是只有84分，现在还在为那一学期的绩点摆烂而还债TT）。}
\SourceParagraph{这边为会这篇材料会先讲一下南京大学的留学现状。详细解读留学要什么申请材料，每个材料有什么用，每个材料的重要性优先级，怎么提升自己各个材料的竞争力。最后讲一下如何选校和学校排名。主要还是从美研研究型硕士甚至是美博的角度开展开（因为要求难度高，但是可以向下兼容）。如果准备英港新思路也适用，如果有申请美研的条件去申请英港新基本上是降维打击。（如果只考虑英港新硕士没有这么难，稳住GPA就行，别的都是锦上添花）}
\subsubsection{一、留学现状}
\SourceParagraph{注：这小一部分十分冗长，如果你不想深入了解留学可以不看。}
\SourceBlankLine[1]
\SourceParagraph{统计大气学院的就业报告，从25届到22届，可以确定的本科出国人数至少为11人、6人、5人、5人，具体去处如下：}
\begin{center}
\small
\setlength{\tabcolsep}{2.5pt}
\renewcommand{\arraystretch}{1.22}
\begin{longtable}{|>{\RaggedRight\arraybackslash}p{0.28\textwidth}|>{\RaggedRight\arraybackslash}p{0.61\textwidth}|}
\hline
毕业时间 & 去处 \\
\hline
2025届、2024届 & 康奈尔大学 \\
\hline
2025届 & 约翰霍普金斯大学 \\
\hline
2025届 & 香港大学 \\
\hline
2025届、2022届 & 新加坡国立大学 \\
\hline
2025届 & 香港理工大学 \\
\hline
2025届 & 墨尔本大学 \\
\hline
2025届、2020届 & 伯明翰大学 \\
\hline
2025届、2022届 & 香港城市大学 \\
\hline
2025届 & 赫尔辛基大学 \\
\hline
2025届 & 北卡罗纳大学教堂山分校 \\
\hline
2025届 & 澳门科技大学 \\
\hline
2024届（两人） & 纽约大学 \\
\hline
2024届、2022届 & 香港科技大学 \\
\hline
2024届 & 雷丁大学 \\
\hline
2024届 & 慕尼黑工业大学 \\
\hline
2024届 & 澳大利亚国立大学 \\
\hline
2023届 & 华盛顿大学 \\
\hline
2023届、2022届 & 加利福尼亚大学洛杉矶分校 \\
\hline
2023届 & 伊利诺伊大学厄巴纳-香槟分校 \\
\hline
2023届 & 匹兹堡大学 \\
\hline
2023届 & 美国东北大学 \\
\hline
2022届 & 美国艺术中心设计学院 \\
\hline
\end{longtable}
\end{center}
\SourceBlankLine[1]
\SourceParagraph{历史上（17届至21届）较为突出且人数多于2人的出国去向还包括普林斯顿大学（17、18届各一人）、苏黎世联邦理工学院（17、20届各一人）、卡内基梅隆大学（17届一人、19届两人）、密歇根大学安娜堡分校（18届一人、19届三人）、曼彻斯特大学（17、20届各一人）。以上去处并非都是在大气专业深造，例如25届去约翰霍普金斯大学及北卡罗纳大学教堂山分校的两人均为转专业，但由于笔者并没有掌握所有转专业毕业生的信息，考虑到删除或标注不全面的信息可能会产生误导，所以上述信息在呈现时没有经过过多修饰。}
\SourceParagraph{刚入校时，笔者曾去新x方、x航道咨询过留学申请相关的事宜（不推荐太早找中介，特别是x方x道）。“大气科学作为一个冷门专业，出国容易申请到好学校。”我偶尔会听到中介这么说。在统计的意义上，大气专业申请留学的门槛相较于热门专业确实会更低。可另一方面，好与坏的定义尚且不分明，上面列出的诸多去向到底是好还是坏，相信每个人的判断都有所不同。我这里客观的理性的来分析一下现状和原因。相信在深刻理解后，你一定能做到理性认知，不再气馁。}
\SourceListItem{1. 南京大学的留学困境}
\SourceParagraph{南京大学虽然是C7，但是在留学方面输了前面几个学校。第一学校层次，不可否认，南京大学绝对是一流大学，但是就大气科学留学来说，南大的头上至少还有北大压着，而且人家北大是物理系的去申请大气科学，竞争力真的没法比。第二留学氛围，其实抛开南赫以外，南大的留学氛围没有清北复交浙浓烈，南京大学的学生真的如校训一样“诚朴雄伟”。（我现在逐渐怀疑南赫的留学氛围似乎也不是十分浓厚，很多人是准备保研的。造成NJU的大伙不留学的原因有很多，包扩理念认同，职业发展，经济情况，校外公司宣传，留学资源等诸多方面）。这就导致我们整体的出国人数不多，且顶尖的苗子未必出国，甚至有人拿到了好大学的offer因为费用问题没有就读，数据自然差强人意。第三留学校外资源，北京上海浙江有很多留学大公司和私人工作室，经验丰富，南京作为一线城市在研究生申请产业上就略显幽默了，毕竟南京的大学中里能去往好学校留学的只有NJU了，总的市场相对来说还是比较低端（在这里就不再更多拉踩别的学校了）；南京确实也很国际化，但是主要是南京的中学生上私立中学申请海外本科（本科和硕士的申请差别很大）；第四校内资源，其实留学的大盘是很看校内资源的，为什么一些传统的中外合办学校，如宁诺，西浦，上纽，昆杜留学大盘那么好，因为他们校内做科研，拿推荐信，成绩单等各个方面都占优势。而在南京大学你这些都要自行安排，磨练修行。资源方面最重要的是connection，清北复交的同学教授大部分都有海外背景，可以教授写强推和并且更易结识大佬教授，所以出国留学优势很大。南京大学因为一些历史原因现在海外资源与清北华五有一定差距，可以说是“嚼得菜根”了。}
\SourceListItem{2. 数据可能并没有很强的参考性}
\SourceParagraph{被哈耶普斯麻牛剑等超级神校录取的人本身就很少，这些学校的有些项目录取率只有2-5\%，而且人家也不止录取中国学生，可能最后一个项目一年只录取一两个陆本学生（传统中外合作办学不算陆本，因为他们可以用合作院校的学历申请，要当作海本看待。而南赫是否能有这点优势存疑）。这个数据量参考意义已经不大，更何况我们这里的高素质选手未必申请出国。其次是理工科的录取难度很大，如果你是文科去一个冷门专业，难度真的降低。同时一些中介给出的留学数据也有造假，这在第三篇中介篇会提及。最后我们深入思考一件问题，我就算今天告诉你南京大学录取哈耶普斯麻的比例是99\%，你就可以躺平了吗，你就不会是那1\%吗；就算我们录取哈耶普斯麻往年的概率是1\%，你就不能成为那1\%吗。与其在这里纠结，不如看两篇文献（但是关注往届的留学数据是有指导意义和反向规划的参考性的）。}
\SourceListItem{3. 留学准备的不充分}
\SourceParagraph{很多去留学的人可能自己当初准备的时候就没有准备好。比如暑研，套辞，推荐信，打磨文书，提前规划，这些费时费力还有信息差的东西真的有那么多人做到了吗。虽然难以估量具体有多少人留学没准备到位，但肯定有人是被信息差或者中介搞了。}
\SourceListItem{4. 南赫学院具有留学优势}
\SourceParagraph{南赫学院的留学优势首先会想到的是外方文凭，但是很残酷的事告诉你，我个人主观觉得这不是主要优势，除非你准备申请赫大和欧洲学校。我们的海本优势主要是指英美港新本科，尚不明确赫大在不在pool和list中（pool是指竞争池，比如陆本在一个pool，海本在一个pool，两个pool中间并不构成竞争；list是指很多学校在判断留学生时会有一份学校清单，在这个清单上的学校会特别关注，甚至gpa要求可能都不一样，比如一个南京大学的均分88和一个普通211的均分95不是一个概念，而不在list中的学校甚至可能直接不考虑录取）。我们南赫真正优势在于gpa友善（你们自己就读了就可以体会到，给分比别的专业好），外语教学（外语科研环境是一个能力问题，留学申请时很关键），科研经历多（暑假有去赫大的科考，并且南赫老师多，方便进组科研），老师有海外经历（你可以去和那些有海外求学经历的老师交流留学经验，并且这些老师写的推荐信作用一定会很大，因为他们有海外的connection）。综上我觉得南赫的留学前景还是会比较明朗的，再说了别的学院也没有我们这么好的氛围，学长学姐们真的是为你们掏心掏肺编写的“难喝生存手册”QwQ。}
\SourceListItem{5. 重新捡回留学的初心}
\SourceParagraph{其实如果真的只是想要做科研，学校的层次没必要那么高，也没必要执着于名}
\SourceParagraph{校，美国前30的博士都是很不错的（后面讲选校时会细致分析），不用怕没学上。如果你只是为了镀金，那么那些水硕交钱又何妨。来了南赫，学科质量（大气科学A+），家庭背景（你都考虑留学了，我默认你家庭背景还不错），和友善的学长/学姐/教授，将会是你的优势。让留学回归本质，只是去境外学习，放下心中的虚荣、攀比和焦虑。}
\SourceBlankLine[1]
\SourceParagraph{接下来就是大家最关注的留学申请材料了。申请制学校有很大的不确定性，没有说达到什么样子就一定录取，或者说一定不录取。我这边是给出一个大致的建议和介绍。}
\subsubsection{二、绩点——究竟重不重要}
\SourceParagraph{虽然本手册一直在劝说大家不要做功绩主义的走狗，让大家摆脱高中思维。但是请记住，这是一个社达的世界（社会达尔文主义），而且对于选拔制，对于普通人来说最好的就是衡水模式（这个模式确实很畸形，但是我们目前无法改变，只能顺从）。留学，去好学校留学，去有含金量的学校留学，注定是一场艰辛的道路}
\SourceParagraph{网上关于留学绩点重不重要这件事有很多争论。真实情况是这样：连绩点这么简单的东西都做不好，招生官凭什么相信你别的事情有能力做好。且申请越好的学校，绩点越重要。绩点的好是没有上限的，90，92，95都是有本质差别的。GPA是绝对可以体现一个人的综合能力至少是自学能力的，他是老师了解你的第一印象。}
\SourceParagraph{首先讲一下留学的绩点和我们保研的绩点处理是不同的，并且你去不同国家和项目的GPA处理和参考程度也是不同的，不同层次的本科院校的GPA也不一样（比如清华大学，南京大学，苏州大学的90分含金量都不一样）。}
\SourceParagraph{留学的GPA我们一般是看全学分绩点，就是说所有的科目都要好好学。而且我们的GPA在申请的时候是只看前三年（申请一般是在大四上，这个时候大四的成绩都没有，但是之后入学的时候要补交大四的成绩）。同时我们的绩点大概率是要通过换算的，常见的是通过WES（World Education Services）换算，这个是将你的成绩转化成北美的成绩单，是申请北美的研究生一定要转化的。}
\SourceLabel{根据留学GPA计算的大致原理，我们可以给出以下骚操作：}
\SourceListItem{1、你可以故意选一些给分很好的通识课拉高你的绩点。}
\SourceListItem{2、你可以把一些不太重要的课推到大四再上，这样前三年卷绩的压力会小一点，大四你只要保证那几门课别把对你的绩点造成毁灭性打击或者挂科，录取的学校是不太可能退你的录取。}
\SourceListItem{3、WES的转化是有一些好处的，比如说90分转化到美国是A，专门课的成绩就是4.0，等于你这门课是满分。但是89分就会滑档，只有3.7。且国外大学又一些课是对应不了的，不如说马原和军理，英语课，可能转化成WES后这门课的学分就会很低甚至没有（请具体看WES官网的转化细则，我这介绍可能有错误，据说不同大学进行WES转化还不一样）。通常来说转化成WES你的绩点会变高。}
\SourceListItem{4、灵活利用学校政策：选课，退课，缓考，红黑榜。}
\SourceListItem{5、主动一点，找老师聊聊学习，学长学姐聊一聊方法。}
\SourceListItem{6、一些阴招，比如说和老师拉好关系，爆金币让同学或者学长学姐辅导你，或者校外机构补课。}
\SourceParagraph{GPA的参考程度在不同的项目和地区也是不一样的，这个你要自己去打听你想上的项目往年申请结果，我这里大致给出：}
\SourceParagraph{英体制（英港新）硕士：GPA占比最大}
\SourceParagraph{美体系硕士：GPA = 论文 = 推荐信}
\SourceParagraph{博士和部分强校研究型硕士：推荐信 > 论文 > GPA（GPA是基本）}
\SourceBlankLine[1]
\SourceParagraph{不要迷信中介给出的低绩点录取好学校案例，因为这是竞争性真相。如果一个学生绩点不出色还能录取好学校一定是因为他有别的过人之处打动的招生官，而这个过人的亮点是很难复刻的。屁股决定脑袋，如果中介不这么宣传，它还怎么有业绩啊。}
\SourceParagraph{目前我知道的只有一种情况你可以不在乎绩点，那就是你的科研十分出彩且申请博士（一定要是博士申请，硕士申请时很在乎GPA的，而博士更看重的是科研能力）。你要出彩到本科时候就有顶刊。但是这一点一般只有计算机好实现（特别是近年来的CV计算机视觉方向），而且科研的不确定性比卷绩点大得多，你所有时间都要all in在实验室里。同时你的绩点不能过低，科研和绩点只能一定程度上置换，本质上是让招生官认为你沉迷于科研而不能兼顾绩点，解释为何你的绩点不够突出（因为这类人太少了我很难找到确切的南大样本举例，但是确实存在比如B站up“whynotTV”的GPA就不是高的吓人，但人家本科发了顶刊还是上海交通大学的计算机，被CMU/MIT/Harvard的计算机/机器人录取了。这个难度在选校介绍中你会了解到）}
\SourceBlankLine[1]
\SourceParagraph{请了解你想上的学校和项目GPA大致需要多少并且为之努力，越高越好，这边给出大致的中位数的参考：}
\SourceParagraph{哈耶普斯麻牛剑\&一些排名世界第一的热门专业（CMU的计算机和机器人）硕士以及其博士， 不低于3.85/4.00 ，最好是3.90+至4.00。但是记住，这些学校是不可能只因为你GPA高而录取你的，但是你GPA低肯定是不鸟你的，除非你是上面说的一种特殊情况。但往往能申请到这类学校的大神们就是全面开花。}
\SourceParagraph{港三新二超热门专业如AI硕士，对口专业GPA系内top3 or 有顶刊。英体制的学校真的有可能因为GPA很高录取你}
\SourceParagraph{IC/港二新三一般专业硕士，GPA3.85/4.00只要能有个正常的科研背景，走正常的机构规划并且完成都能大概率被录取。}
\SourceBlankLine[1]
\SourceLabel{但是大家不要有卷绩焦虑：}
\SourceListItem{1.绩点比你高的同学可能和你不构成竞争。比如说你虽然是3.85/4.00，原始分90/100，在你们这一届排名只能排30\%（说的就是我们这一届），但是你前面的同学有没有可能科研完全不如你或者说人家就没准备出国}
\SourceListItem{2.除了绩点是成绩参考外，还有一个重要的东西是排名。排名是否提交首先要看这个申请的学校是否需要你提供排名且你交的成绩单中是否包含排名。排名的作用主要是在你绩点低但是排名高的时候补充说明，比如说你们这一届给分严，你只有3.6/4.00，但是你在你们这一届排名是第一，那就这个成绩就完全没有问题，你就可以在提交绩点以外额外提交排名。或者说你们这一届全部绩点都很高，你3.85/4.00都不到前10\%，那你就不要提交排名。}
\SourceListItem{3.绩点上的劣势还是能通过别的方面来尽量置换弥补的，比如说科研十分优秀，推荐信件十分优秀，文书十分有个人特色。}
\SourceListItem{4.你的GPA的构成也会让你的GPA变好看，比如说你的GPA逐年变高，这样老师会觉得你是一个有潜力的人，只不过一开始不适应大学生活；或者你的专业课GPA都非常高，你的GPA是被那些不咋重要的课程拉低的；或者说你学了一些本硕贯通的硬核课程分数很高}
\SourceBlankLine[1]
\SourceParagraph{但是总的来说，GPA是你的基本生命线，而且是唯一只可以靠你的个人努力改变和决定的东西。}
\subsubsection{三、语言标化——多高才算高}
\SourceParagraph{语言标化是来证明你的语言能力的。主流的是雅思，托佛，GRE，GMAT。这些网上很容易就可以搜到统一公开的资料，不做十分详细的介绍和经验分享。}
\SourceParagraph{雅思是英国学校喜欢用的，但是雅思的特点是低分好拿，高分难刷。}
\SourceParagraph{托佛是美国学校喜欢用的，高分段比雅思好拿。有人说要先托佛后GRE，也有人说要先GRE后托佛，请自己定夺。}
\SourceParagraph{GRE目前有些学校是optional的，也就是可以不提交，因为GRE是体现你的英语学术水平，如果你有英文的论文出版和一篇好的自写文书，GRE真的不重要，可以不考。（但请看看清楚你想去的项目要不要GRE）}
\SourceParagraph{GMAT是商科用的，不介绍。}
\SourceParagraph{标化的整体准备思路是达线就行，比如哈耶普斯麻建议你托佛115就足够了。}
\SourceParagraph{如果说英体制的学校还有可能因为你GPA高录取你，那么绝对没有一个学校会因为你标化高录取你。标化只是一个招生筛选，达线即可，请不要过度投入，它绝对不是重点。}
\SourceParagraph{标化的准备是最需要使用互联网的，这玩意基本公开透明资料多，可以自行准备，实在不行可以报班上，而且这玩意可以氪金（这方面100\%黑产，但是我在下一篇中介那边会讲。千万不要弄虚作假）}
\SourceParagraph{标化还值得注意的是它的年限。托佛有效期只有两年。并且标化不仅申请时有用，在申请暑期科研的时候往往也要标化。}
\subsubsection{四、文书——最容易被忽略的东西}
\SourceParagraph{文书本应当，也应当是你自己真实的创作，因为这是招生官了解你的最好途径。一切都是你真实的信息。在招生中特别是硕士招生中十分重要，等于说是你的文书代替你去招生官面前面试。}
\SourceParagraph{请警惕一些中介和你说的文书不重要，因为大部分中介没有能力给你处理好文书！所以它只能这么说。而且一个标准化，流水线，没有特色，且可能仅是一个英语专八的文书老师给你写的文书，在留学申请中的作用也堪称是不重要了。}
\SourceParagraph{文书中的门道很多，原谅我现在还没到申请季，所以没有了解特别多。以后会补充，你可以先去B站搜集一下信息，强烈推荐B站up主“泡芙方”的文书写作课指导，免费的，我只看了一两集，但是他从一个一线美国教授的角度给出了很多很有用的建议（此人后面有详细介绍）。}
\SourceParagraph{总的来说文书要提供个人信息，力求真实（可以利用叙述包装，但是别弄虚作假），要能有个人特色使你脱颖而出，塑造西方精英形象，符合西方写作习惯以体现你英语能力。博士的文书要体现出你强烈的科研愿景和热情，因为教授是来培养你做科研的。}
\SourceParagraph{文书分为CV，SOP。而且硕士的文书和博士的文书有大不同，不同学校不同项目对文书的要求也不同。请具体情况具体分析，按照你想录取的项目的要求走。}
\SourceParagraph{文书是一定要写多版不断打磨的，这期间你可以多使用一些工具，如AI（别全用），高手修改（有专门的文书网站，外国人一线人员亲自修改，次数收费），机构或中介修改/润色（这个会在之后的中介篇写）。给不同的学校也可以针对性的写不同文书。}
\subsubsection{五、对口程度——申请的前提}
\SourceParagraph{我们大气主要对口的是：1大气与海洋科学学校2土木与环境工程学院3地球与行星科学学院等地科4环境学院。当然也有稍微偏一点的数据科学（是计算机大类），金融（大气环境金融衍生类产品），但是计算机和商科是出了名的卷，你专业还不是完全对口，你申请就要接受学校档次可能要向下调整。}
\SourceParagraph{你真的想换赛道的话还是有一些技巧可以用到的，比如说免修不免考来选更多的课和跨校区选课，用coursera补充海外课程，在科研中使用横向交叉领域论题（南大真的有通过AI+大气的科研论文最后成功转AI的例子）。就算不能直接申请到别的专业，国外的制度和专业也是很神奇的，比如你可以在斯坦福的土木与环境工程学院学很多计算机、AI、工程相关。}
\SourceBlankLine[1]
\subsubsection{六、活动/实习/科研——软背景其实很hard（硬）}
\SourceParagraph{我这边主要讲科研，理工科主要看科研。科研水平是一个教授是否收取你当博士生的重要参考，对于硕士来说也是能使你脱颖而出的加分项。}
\SourceParagraph{首先明确，大部分的科研不需要天才，特别是engineering项目，是看你的投入有多少。除非你是搞理论研究，这个玩意太吃天赋了，而且这个玩意不需要导师指导。}
\SourceParagraph{科研你最好要热爱再去做，这样你才能用你的休息时间以及寒暑假时间去做。如果你不热爱，只是为了出国科研，说实话有点本末倒置。但是你这么做一定有你的原因，至少选一个你不讨厌的方向。}
\SourceParagraph{总的来说你要比组里master还有phD做的更好才能做出出色的科研。}
\SourceParagraph{要做科研的话就要进组。请不要羞涩！主动约老师谈，发邮件，找office hour，且这也是一个和老师建立connection的过程。本科生进组老师一般来说是很欢迎的。而且你完全可以表露你想出国的心愿。但是尽量大二再进组，因为老师们平时都在苏州，且大一你啥都不会。顺带一提可以跨院系跨校区找老师做科研，只要老师愿意收你。}
\SourceParagraph{一开始进组肯定是啥也不会的，可以多找学长学姐学习。请不要急于求成，一开始都是从读文献做起，然后普遍大二大三才能有进展。}
\SourceLabel{科研设计到几个小问题这边总结一下：}
\SourceListItem{1. AP组 和大佬组。AP组指的是刚从国外回来任教职的老师，一般来说还有一线科研的能力，思维活跃，点子多，愿意培养新人，但组内人少，要主动试错。而大佬组就是那种帽子特别大的老板，他们名声大，好发文章，发文章远远快于AP组，总资源多，但是问题是组太大可能要博士生来带你，还不容易拿到牛推（这个在推荐信说，就是说推荐人是牛人，但是因为和你不太熟不一定给牛推）。这个问题没有答案，和找对象一样，适合自己的最好。有关导师是怎么样的可以去看B站毕导的相关介绍《【毕导】从博士到教授到院士，中国科研的打怪升级之路》}
\begin{figure}[H]
  \centering
  \includegraphics[width=0.92\textwidth,height=0.78\textheight,keepaspectratio]{figures/wq\_image3.png}
\end{figure}
\SourceListItem{2. 三个月还是三年。一个组待三年还说三个月就换组，这个东西在于多尝试，并且适合自己的策略。而且要有第一段打杂的经历才能有第二段承担关键环节的能力}
\SourceListItem{3. 1作还是n作。这个东西很难直接说，这个涉及到你发的是什么期刊，你目前有多少期刊}
\SourceListItem{4. 没有产出？产出我们一般指的是论文，我明确的说，没有产出是可能可以申请到好大学的！因为大学的招生老师和博士可以通过你的SOP和你的面试知道你有没有科研能力和科研思维。但是有产出是绝对好过没产出的。}
\SourceListItem{5. 3+1\&暑研\&国际会议。这个十分重要，就是指你有没有海外经历，3+1是指你大学有一年交换去国外，在那边边修满GPA边找好的课题组做科研。暑研是指你大三升大四的暑假套辞到国外教授去现在做实验。国际会议在大气方面只要是指AGU的年会。这方面涉及到你的外方connection，外方强推牛推，外方科研产出，外方教授结识，我会在推荐信那一块再次提及。我可以给你保证，所有能申请到哈耶普斯麻层次学校的学生，都有很好的外方connection。}
\SourceListItem{6. 英语汇报能力，包装能力。组会汇报多练习，尽量用英语汇报。英语学术能力很重要，这个是硬能力。}
\subsubsection{七、推荐信\&connection——半只脚踏入梦校}
\SourceParagraph{选择暑研\&3+1\&国际会议基本上就要因为时间冲突放弃保研夏立营了。参加这些项目最重要的就是就是科研、推荐信、connection。科研之前说过了，这边主要讲推荐信和connection。特别是对于博士申请来说，基本就看三个东西：你的本科，你的成绩，你的connection。}
\SourceParagraph{暑研基本上要从6月到11月。这个是你要自己找教授套辞发邮件，然后得到同意认可再去现在做的，而且基本上全程自费。等我以后申请暑研会再补充上来，这个玩意是要DIY的。3+1我们学校有部分学校的合作，可以在群内查找，你在交换期间要边修满绩边找导师做科研。参加学术会议主要目的是去认识大佬，和大佬交流，如果顺利可以得到大佬的口头承诺入组当博士生。}
\SourceParagraph{这些项目可以说是你的博士生demo版，你这样就可以提前了解到国外课题组的氛围，提前了解过外国的国情，了解外国路线。}
\SourceParagraph{最重要的是，这玩意对申请超级有用。首先就是关系网，科研圈是一个巨大的派系关系网，你要是能在大佬的门下有这个关系网，在以后的申请，面试，套辞时就可以扯上关系，招生官和教授就会对你印象深刻并且额外注意。并且师出同门就意味着你以后可以找到同门的别的教授课题研究，课题研究思路大致是一脉相传的。PH.D.尤其注意，在外方课题组如果你干的好，说不定可以直接PH.D.去往你的暑研课题组。因为PH.D的决定权在教授。我再次强调connection的重要性（connection你可以理解为有在目标学校有影响力的教授强推，或者认识招生的老师和教授），硕士和博士是申请制度，有connection的学生等于有一份信誉保障，招生部和教授对你知根知底，也跟放心录取。（真实案例，作者有认识北大的学生就是因为某个教授每次来国内开讲座都去前排听讲交流，令教授印象深刻，被教授收为博士生）}
\SourceParagraph{其次是推荐信，这是外界对你的客观评价。总的来说，海外的推荐信含金量比国内高，国内有海外背景的老师推荐信比不同老师好。推荐信的含金量主要是看：1你的推荐人是谁2你的推荐人是否和你申请的目标院校有connection和话语权3你的推荐人是否真诚推荐你且给出高度评价。根据这些加权就可以得出一份推荐信是否为强推，例如一份新晋教授的真心推荐信可能比一位业界大牛的流程化推荐信更有用。因为国外是有一种信任体制的，写一份推荐信是要担责的，这代表这个推荐人的选人眼光，所以申请时对于海外推荐信的认可度十分高。如果你能拿到一份你暑研期间/交换期间/学术会议期间的老板的强推，你就获得了当地的connection和绝对的亮点。但是就是说你要猛猛干，势必拿到强推和产出，否则暑研\&3+1\&学术会议都会血亏。（这些都是要自己花钱的）}
\subsubsection{八、什么申请材料没用}
\SourceParagraph{有些材料的作用真的聊胜于无，甚至可能是扣分项，他们或多或少都沾一以下几点：无特色，大众化，知名度低，门槛低。}
\SourceParagraph{首先就是各类国内竞赛，作者一开始大学是想搞ACM的，这个东西要是能拿到金牌是一定可以去清北的，但是很不幸，大多国外教授是没听说过的。同理，各类大创，省赛，国奖可能都没有用。但是还是有一些竞赛是有点用的，他们的特点是全球比赛且知名度大，比如说美赛（数学建模），Kaggle（数据科学比赛），甚至github上开源项目（虽然这不算比赛）。}
\SourceParagraph{其次是你各种能在低质量中介上能买到的便宜科研和推荐信。大致判断标准是：1你在各个平台上都看到过（我真的有在两个不同的中介那边看到过同一个科研项目）。2是一个多人线上科研。3价格不算离谱（真正的私人定制科研是要6、7位数的，但是你们也千万别被骗了，而且这个玩意是有点灰产的，后面一篇中介那边会细说）。总结来说就是这些东西是那些中介既可以卖给A同时又可以卖给B的，是偏向大众化的。这些大多没啥用。}
\SourceParagraph{还有各类没啥个人能力的门槛的，这边主要指的是一些暑校。这玩意的门槛是真的低，点击即送，主要作用就是感受大学氛围。真正有含金量的是暑研，要套辞还要去线下教授实验室，混的好可以有科研产出和强推。}
\SourceParagraph{总的来说真正有用的申请材料要么是靠个人能力挣来的或者家族势力托举来的，且有自我特色的，这些才能让你在一众申请者中脱颖而出。留学终究是一场对综合能力的考验（这边稍微讲一下家族托举，首先我觉得一个人合理合法合规的利用周围条件达到目的是没用问题的。自己家里有能量且可以调动为什么不算能力的一种？这也就是为什么藤校偏向于录取有家族背景的同学。家里的托举是要有家里有很大能量的，比如说有私交非常好的领域大牛写的推荐信，家里是学术世家可以在科研上提供帮助甚至是找人指导，家族企业可以提供优质实习岗位，家里有资本可以为孩子先搞到科技专利再为其融资，家里资金和门路将钱置换成上类讲的各种资源。家境一般的同学也不要自卑气馁，人只能和自己比较，相信自己可以用双拳打破家徒四壁）。}
\SourceBlankLine[1]
\SourceParagraph{申请材料总结，总的来说申请材料就是给招生官展示出一个完整立体的你。本质上来说反应的是你的能力，但是也要平时也要额外注意积累，很多是没法临时抱佛脚的，这么多材料交上去加上可能存在的面试让你无法隐藏（除非你用大资源给自己搭建一个人设，补足短板）。准备和积累申请材料的过程本质上是一条锻炼和增长能力的一条路。}
\SourceBlankLine[1]
\SourceParagraph{总的申请时间线（来自于B站up“博可菌”，此人上海科技大学申请到了加州理工的博士）}
\begin{figure}[H]
  \centering
  \includegraphics[width=0.92\textwidth,height=0.78\textheight,keepaspectratio]{figures/wq\_image4.png}
\end{figure}
\SourceParagraph{南赫大一学分占比超级大，好好卷GPA。请好好使用大一，这是你最闲的时候}
\SourceParagraph{大二我还没经历，目前是准备边搞科研边卷GPA。}
\SourceParagraph{大二升大三的暑假目前明确知道有暑期学校，可能暑假会被压缩。}
\subsubsection{九、学校和排名}
\SourceParagraph{排名上主要是有qs排名，USsnews排名和软科。里面又分为综合排名和专业排名。}
\SourceParagraph{qs排名是英国搞的，逐渐不靠谱，偏向英国体制学校，是一个充钱就可以变强的榜单（排在第二的IC，第九的NUS/UCL，比yale高一名的新南威尔士，比清北高的港大，比南大高的港理工，年年升高的澳八大都是名场面）。}
\SourceParagraph{对于不懂留学的人和一些公司的HR，他们在看留学生时只在乎qs排名。qs排名虽然不精准但是普及，名声大。再者很多人都是这个机制的受益者，比如一些家庭有能力的人但是孩子一般的人会把孩子送往澳洲留学，澳八大的整体qs都很高，所以qs排名还不能弃用。}
\SourceParagraph{USnews是美国弄的，显然更加合理，但是不可避免的偏向美国学校，而且有一些学校公立大学的评价似乎过于高的（比如UW，UCLA，UCSD）。但是还是比qs排名合理，去美国留学可以主要参考这个榜单。同时USnews是有专业榜单的，这个十分重要。比如说一些从事科研的人，是完全可以按照专业榜单去申请而不看学校整体titile的，比如UW的大气科学是世界第一，CMU的计算机是世界第一。美国的圈子还是非常看重专业榜单的，和国内不太一样。}
\SourceParagraph{软科排名是国内榜单，近几年逐渐出圈。里面的综合排名和专业排名是非常符合国人价值观的，而且基本只看学术能力，没有一些奇怪的指标。推荐参考，但是它毕竟不是留学榜单，还是要辅助别的榜单参考的。}
\SourceParagraph{所有的排名只能参考，主要作用还是让你知道有哪些学校，混个眼熟，深入了解才能发现里面的大学问，什么互联网学历鄙视圈，录取难度，毕业年薪，王牌专业，校友资源，地理位置，地区认可程度。不可迷信排名，请根据自己情况细致了解。到这种程度的问题就可以去reddit（约等于外国的知乎）和一亩三分地（国内最专业的北美留学平台）上仔细询问和查看。选校时请以一手资料和学校官网资料为第一指导！}
\SourceBlankLine[1]
\SourceParagraph{总体选校的时候要综合考虑：适配度、城市、排名、个人发展、家庭条件、自身能否申请上。首先明确一点，我们虽然说是在选校，但是本质上是在选项目。相同的学校不同的项目是不同的，比如说宾大的硕士和博士不是一个bar（难度和含金量上宾大硕士是显著低于宾大博士的）}
\SourceParagraph{总之在选校时有不可能三角：含金量高/门槛低/花销低。}
\SourceParagraph{当然有含金量高，难度大，开销大的，比如你想要硕士去读哈耶普斯麻。}
\SourceParagraph{总的来说这三个东西也是可以互相置换的，比如你要是家里有资源可以给你安排活动、科研、补课，那么申请难度就会稍微小一点。}
\SourceParagraph{选校是一件很个性化的事情，主要是要从专业，学校title，学校地理位置考虑。这三点是决定你的发展的（比如纽约大学的地理位置是极其好的，虽然他只是美国top30的学校）。然后根据自己的能力给出一个合理的选校单，进行逆向规划。选校理论上是不设置上限的，所以说和高考志愿一样要有冲，稳，保三个档次，比高考志愿好的点在于可以拿到录取后自己抉择去哪一所学校。}
\SourceParagraph{选校主要是为了有大致的目标和方便逆向规划，比如找到你想申请的学校的项目的往届录取学生数据，以此为目标努力。但是在申请季前的选校没有太大意义，在申请之间，你没有必要考虑你是更适合哈佛还是更适合耶鲁。}
\SourceBlankLine[1]
\SourceParagraph{在第一版中本人列了一个留学学校从夯到拉。但是现在回头看感觉怪怪的。总的来说可以按照录取率和毕业学生去向来判断。基本上可以参考一些顶级公司都收那里的毕业生以及其数量来参考（比如说可以参照一些硅谷顶级AI科技公司，或者华尔街顶级量化与对冲基金如：Openai，Jane Street，BlackStone ，JPMorgan这些）。同时提醒一句，有一些硕士项目，光是南大本科的bar是很难够得上的（要叠加顶会论文，北美大牛推荐信），几年才能出一个。如果非要冲那些大藤项目（例如普林的Mfin），可以考虑转学。（转学以后可能会专门开一章节来说）}
\SourceBlankLine[1]
\SourceParagraph{下一期会着重介绍有关中介和DIY的内容，会涉及到灰产和留学公司的恶心操作。}
\SourceBlankLine[1]
\subsection{留学3——DIY和找中介QwQ}
\subsubsection{一、中介的作用}
\SourceParagraph{说实话仔细想想中介的存在有一点不合理，因为国外的学生自己申请本国的学校是大概率不会找中介的。}
\SourceParagraph{其实最初来说中介就知识提供信息源，帮助留学季，一般来说分为半包和全包，期中最推荐半包，半包适合自主性强和申请博士的同学。全包就是说包括：培养方案建议，文书写作（可以本人不参与，但是不建议文书不参与，我已经说过文书的重要性），选校，头脑风暴（一般来说是介绍专业，帮助选校等一系列），模拟面试，帮助提交申请资料，帮助拿护照，极端情况下的全包你是可以什么都不干的。半包就是基本只提供信息，提供选校，给点建议并且修改文书。现在中介都很卷，所以基本不管是全包还是全包都会给你找海外对口专业的学长学姐提供建议。}
\SourceParagraph{但近年来中介也开始提供一些新的业务，比如提前规划和背景提升。但是说实话，这些东西含金量低且水深，而且有校内资源的话显然更好。基本上没有人光通过中介就逆天改命的}
\SourceParagraph{最后还有一些一般中介不做的灰产服务，我之后单独说明。不要碰！不要碰！不要碰！}
\SourceParagraph{当然有部分中介真的能做到一些事情，比如说工作内推，校友内推，私人科研团队指导。但这玩意一来搬不上台面，二来要氪金（大几万到几十万不等）。}
\subsubsection{二、中介里面的恶心事}
\SourceParagraph{跑路。近年来有几个中介因为资金链出现问题直接倒闭跑路了，比较出名的是再来人。而且近两年留学行情大盘越来越淡，很多种中小型公司都会跑路。而一些小型公司可能还好一点，如果老板自己是法人且有留学经验，可能会把你这一单自己来接手做完。但是总之跑路问题确实一直存在。}
\SourceParagraph{恶意引流。小红书上很多帖子，群聊甚至声称做机构避雷的人都是托。某些的大机构一年在宣传上花费的钱就有7，8位数。而那些说自己是避雷博主的人更加是脸都不要了，给你避雷拉踩完后给你推荐机构，其实就是收了钱的托。}
\SourceParagraph{数据造假。有的机构会找到一些根本不是自己做的留学数据进行造假。例如你是一个海本申请到了前五的学校，中介联系你给你塞钱、让你说你是在这家中介做的，所以你能看到有一些中介的哈耶普斯麻申请数据好到不正常。那些申请老师的背景学历可能也是假的。}
\SourceParagraph{信息差。现在很多中介会说申请不上前10的学校就退钱，但是前十的学校有一些是很水的，比如说宾大和哥大的一些专业。}
\SourceParagraph{虚假宣传。和你一开始交流签单子的人是公司里的销售！而负责你的是申请老师。销售说的天花乱坠，但是申请老师可能是个棒槌。还有一些你无法追查和监控的细节会存在虚假宣传，比如申请老师的经验和学历。}
\SourceParagraph{工作质量。一些大机构如新东方一个老师可能要带30个学生，实在是多的过于繁忙，让你DIY都是小事，甚至有的时候耽误你事情。而且老师能力不行，有的文书老师自己可能就是英专八级，而根本没有写文书的经验。}
\SourceParagraph{附加消费。有一些中介会pua你说你软背景不行，然后让你买它们家的科研活动。}
\SourceParagraph{相互拉踩。不用的公司之间相互攻击，影响你的判断。}
\SourceParagraph{合同踩雷。 合同中的退款不合理，且有一些前置要求（要求如果语言，gpa不到要求分数，合同失效）}
\SourceParagraph{请记住，这里面水很深，请仔细甄别，保护好自己。}
\SourceParagraph{还有就是现在留学公司普遍都比较穷，所以质量都不太好。}
\subsubsection{三、究竟要不要请中介}
\SourceParagraph{从上面的介绍你大概就可以知道了，中介这玩意只有不靠谱和十分不靠谱。但是中介这个行业能存在就说明它还是有市场的。}
\SourceParagraph{中介最重要的作用是提供信息，虽然这些信息都是网上公开透明的可查找的。但是大部分人不会魔法也不知道去哪里找，而且中介总归可以省时省力，有更多时间干别的事。}
\SourceParagraph{其次中介还是有一定的督促作用的，找个push一点的中介，对p人来说还是很重要的。}
\SourceParagraph{最后就是中介真的很能提供情绪价值，如果你没有留学搭子，中介是你焦虑时唯一可以倾诉的对象。}
\SourceParagraph{总的来说基本上所有人都会在申请时感慨：好像中介没什么用。但是大家很难摆脱中介，因为大家一直以来出国很少不找中介的，很少有人有那个敢自己走在留学这条路上自己承担责任。}
\SourceParagraph{这边建议找中介时尽量找半包+DIY，特别是申请博士的，因为中介在申请博士的作用会减低很多。中介一般都说自己擅长包装，而博士申请是要面试的，所有的包装都会在你和教授沟通时被拆穿。}
\SourceParagraph{如果真的舍得花钱找中介尽量找有口碑的中小型工作室，因为他们是可以对你的上限提升的，但这些工作室可能抗风险能力不足，而且价格昂贵。大型留学机构一般都是流水线。找中介真的不急，适合自己的是最好的，尽量多找几家聊聊，慢点做决定。注意所有中介负责和你聊天的都是销售，真正重要的是你的申请老师，找的时候最好问这家公司要和你bg匹配的同学的申请情况辅助判断。}
\SourceParagraph{请不要中介所有的话都听，要辩证地看待中介说的话，请进行信息交叉验证，不要所有事情都甩给中介，就算你找了中介，也要留心眼，尽量能自己参与的部分都参与，私下DIY多申请几个学校都行。}
\SourceParagraph{中介可以自行了解，一旦你开始打听，不久你就会被各种中介加满（特别是你开始在小红书上寻找）。}
\SourceParagraph{找中介一定要确定对方的底细，包括给你辅导的老师是什么经验学历。这个机构带出过怎样bg的学生申请上了怎样的学校，公司的法人是谁，退钱怎么办，交流过程能不能录音，营业执照是否正确，合同是否模糊有漏洞。总之一旦发现不对劲，立刻跑路。}
\subsubsection{四、黑产}
\SourceParagraph{这边仅做介绍了解，但是请不要造假和保录取！后果很严重：1你被退学且沦为失信人员。2你会影响学校形象，我们学校已经被MIT金融拉黑，因为有一位申请上岸的同学被发现造假，此后MIT金融不再收南大学生。害人害己。}
\SourceParagraph{保成绩}
\SourceParagraph{保成绩分为语言成绩和学校成绩，但总的来说就是造假。语言成绩的造假方法分为替考和买答案，一般来说比较难查出，但是现在已经越来越困难，因为：1现在大多数学校不接受在家考和去东南亚等别的地区考。而这些考点和方式都是作弊实行的手段。2学校有开学后会自己再测试语言，如果发现水平不一致会严查3现在各个考试换题库平凡，买答案也不一定买中。学校成绩造假常见于高中申请本科，一般来说就是说是中介自己拥有一个院校或者与院校合作，找那个院校给一份假的成绩单。但是现在各大学会找原学校教务处确认成绩，被拆穿的风险很大。还有干脆演都不演给成绩单p图的。。。现在也有那种找人给你下线代修CC课程的，用来CC转UC。}
\SourceParagraph{保录取}
\SourceParagraph{保录取最大的可能是诈骗。常见的诈骗有混淆学校名称和p录取图的，还有直接跑路的。英国和澳洲的学校保录取还是有可能的，特别是ucl，kcl，澳八大（美国的哥大和宾大其实近年来也有保录取丑闻），具体操作我不太清楚，但是总归离不开造假or氪金（这些学校氪金基本要7位数），而且专业一般都很水。美国学校保录取存在，但我们一般来说称之为legacy录取，要么是氪很多金（一般7,8位数以上），要么是你有特殊门道可以找到某个教授or院长，在你特别讨教授or院长喜欢的前提上再小捐一笔。总的来说是常人触碰不到的。。。}
\SourceParagraph{再提醒一遍、这些灰产不要碰！不要碰！不要碰！害人害己！}
\subsubsection{五、DIY申请一些事项}
\SourceParagraph{鸡蛋不要放一个篮子里。现在留学局势变化很大，多申请不同的地区，不同的学校。}
\SourceParagraph{能力至上，总的来说申请留学总归是你的能力问题，你的学习能力，科研能力，展现自己的能力，信息收集的能力，心态调整的能力，执行能力。这些能力的提升是最重要的。}
\SourceParagraph{善于寻求帮助，虽说是DIY，但是我这边还是十分推荐半包的。而且可以单独用一些服务，比如专门找专业人士模拟面试or找专业人士改文书。}
\SourceBlankLine[1]
\subsubsection{六、信息}
\SourceLabel{强推（我这三篇的总结基本都来自于这里）：}
\SourceParagraph{B站 【一天一梦校美国篇】合集！详解美国综合大学top100 （很细致的美国学校讲解）}
\SourceParagraph{B站/youtube“泡芙方ProfFang”（这位是在斯坦佛读的硕士和博士，且在美国匹兹堡大学任教，研究方向为环境流体力学。本科二本转学美国100来名的大学再研究生逆袭斯坦佛的传奇人物。从教授的角度给出了很多建议，做半包和规划服务）}
\SourceParagraph{B站“博可菌”（上科大录取了加州理工大学的博士，在B站和小红书无偿分享自己的留学经验）}
\SourceParagraph{B站“WhynotTV”（上海交大申请到了MIT，哈佛，CMU的博士，在B站无偿分享）}
\SourceParagraph{我强推B站的个人势up主的原因是因为来源可确定，信息包真。}
\SourceParagraph{一亩三分地（\url{https://www.1point3acres.com/}）论坛，北美留学方面很全}
\SourceLabel{其余：}
\SourceParagraph{“BeBeyond”、“棕榈大道”等中介的微信公众号（里面内容都是换汤不换药的，自己找一个看的顺眼的公众号专注一下，当个信息源）}
\SourceParagraph{推荐的信息网站}
\SourceParagraph{推荐网址一，查英国项目：\url{https://www.admissionreport.com/}}
\SourceParagraph{推荐网址二，查全球项目：\url{https://www.thegradcafe.com/}}
\SourceParagraph{The Student Room: \url{https://www.thestudentroom.co.uk/}}
\SourceParagraph{Fiverr 文书润色：\url{https://www.fiverr.com/?source=top_nav}}
\SourceParagraph{open CS App北美申请案例：\url{https://opencs.app/}}
\SourceParagraph{欧陆申请案例：\url{https://global-cs-application.github.io/}}
\SourceParagraph{B站和小红书有很多宝藏，但是记得甄别！那些个人势的，可以查明信息真实性的，没有广告推销的，在一线的（已经在留学或者是教授）博主一般来说都是可信的。而那些上来就让你给电话信息，加个人微信引流的，给你推荐中介的，大概率都是收钱办事的。}
\SourceParagraph{再次重申请不要随便相信各类平台上的（特别是小红书）一系列不清楚对方底细和目的的人/中介，大概率是为了骗你钱来签单子。}
\SourceBlankLine[1]
\SourceParagraph{南赫手册留学篇会持续更新并补充真实案例。}
\section{直接就业}
\SourceParagraph{不建议本科毕业直接就业}
\SourceParagraph{应该不会有人这么打算...吧}
\SourceBlankLine[2]
\section{跨专业}
\SourceParagraph{施工中\textasciitilde{}\textasciitilde{}}
\SourceParagraph{大气跨专业群（QQ）：1057594377}
\chapter{交换学习}

\vspace{-0.5em}\hspace{8em}
{\large\bfseries 当我们谈论交换时，我们在谈论什么}

\vspace{1em}
\SourceParagraph{说明：本文非官方且强个人经验导向。涉及学分认定等制度性内容，均以当年的最新政策为准。文中只描述笔者当时获得的信息，不构成保证。}
\section{Overview}
\SourceParagraph{以下讨论的对象主要是一学期及以上的自费/公费交换生项目。关于 summer school / 暑研 的内容会在文末简单提及，之后也可以单独补充。交换不是一个对所有人都必然更优的选择，但如果你希望暂时离开原有轨道、探索新的课程和方向、尝试科研、体验另一种生活方式，它值得你认真考虑。}
\section{交换学习有哪些好处}
\subsection{1. 交换与绩点及学业}
\SourceParagraph{整体来说，笔者认为交换通常有利于你的绩点。首先，虽然并未在交换结束后转换自己的学分，笔者了解到，交换期间修得的学分（课程）可以自己选择是否转换为南大学分。这意味着你可以只转换得分高的课程，而不转换得分低的课程。利于提升绩点！如果你认真修读一门课程，拿到不错的总评，转换回来可能对成绩单很友好；如果你像笔者一样，交换期间把大量时间投入科研或其他事情，课程成绩一般，也可能存在不转换某些课程的空间。规则允许的前提下，交换往往给了你更大的课程组合与时间安排自由度。}
\SourceParagraph{谈及保研，很多同学会担心，交换期间无法选修部分保研课程，会导致失去保研名额。至少以笔者交换时学院给出的口径来看，交换期间未修得的课程可以在后续学期补修，不会影响到保研名额的分配。某种意义上，这甚至可能形成时间上的利好——你可以把部分比较困难、得分普遍偏低的课程放到后面补修，而把交换学期的时间留给更重要的事情。不过目前南赫第一届学生尚未完整走完相关流程，很多细则仍需要观察。交换期间的选课同样值得认真规划。部分院校会限制交换生只能选择本专业课程，但这类院校并不是多数；更多院校允许交换生跨专业选课，因此对想探索大气交叉方向、AI、CS、遥感、环境、金融等方向，或者想为之后转专业/跨申积累 bg 的同学非常友好。反过来，如果你非常确定自己要继续深耕传统大气，交换院校未必能提供比 NJU/NHI 更完整的大气课程体系，大气毕竟是相对冷门的方向，很多学校相关课程不多，这一点反而可能不如留在 NHI 潜心学习。交换期间选课受限是一个问题，选课优先度又是一个问题：部分院校交换生优先度低于本校学生，部分相持平或更高；有的限制研究生课程/商学院课程，有的则不做限制。我不认为需要过于担心，但在申请前最好自己查清楚。若你投入足够多的时间精力，交换课程还可能带来一两封不错的课程推荐信。比方说 UCB / NUS / HKU 的高阶课程推荐信，相比一封泛泛的 NJU 课程推，对申请境外高校大概更有用；课程推最好来自高阶、能体现能力的课程，类似动力气象，而不是“大气科学概论”这类入门课。}
\SourceParagraph{最后，workload。留在 NHI 的大三上，加上免修不免考的两门课，笔者一共有 12 门期末考试的课程；而在 UST 交换时，项目要求选修 12–16 学分（印象中），笔者总共只上了 4 门课就满足要求。课程负担下降以后，你才真正拥有一整块可以自主支配的时间。交换学期中，你可以重新分配自己的时间，拿更多更自由去科研、实习、读书、社交、旅行，或者只是认真想清楚自己下一步要做什么。以上课程推、学分转换等部分，一些学长学姐曾有过实践；总体来说，笔者认为交换对绩点/学业的影响整体是正面的，制度性内容每年可能会有变化，务必自己确认。}
\subsection{2. 交换与科研}
\SourceParagraph{对于很多研究导向的升学路径而言，绩点很大程度上决定你的下限（能不能进入牌桌），科研经历则决定你的上限（之后能走多远）。科研经历不仅是简历上很重要的一方面，同时也可能带给你 connection / paper / recommendation /reputation。交换期间做科研，我认为有如下好处。首先，是更容易进入一个以产出为导向的科研环境。南大很多老师乐于带本科生体验科研，但整体并没有非常强的“带本科生发论文”风气，更不用说让本科生作为论文的主要完成者、署名一作/共同一作。南大之外，一些高校的评价体系中，发表论文是综测加分中影响很大的版块，综测又与保研资格相关；在这样的指挥棒下，残酷的升学竞争催生了一批本科阶段已经手握多篇论文的申请者。nhi 的同学相较而言有着华五本科的优势，但这往往不能抵消科研经历的差距。面对已有成熟科研经历、手握论文或强推荐信的其他申请者，仅靠学校 title 往往不够。南大并不推崇本科生把发 paper 当作科研的唯一目的，我也认同科研不等于 paper；但在升学游戏中，论文作为一种高度可见、易于识别的科研产出，在申请时手握 n 篇论文往往是优势。境外高校，类似 hk 的一些学校中，很多组的科研更偏论文导向。如果有机会在交换期间找到人带着做科研，这里对你的定位可能不只是“打杂”或者“探索熟悉某个方向”，而是从一开始就围绕一个明确问题、以最后形成完整工作为目标推进。另一方面，华五学生在南大本校当然并不稀缺，但不少境外华人老师平时接触不到那么多华五级别的本科生，所以会愿意接收一个来自南大的学生，尤其是交换生这种能够在线下一起工作一个学期的免费劳动力。如果你展现出了较强的能力、自律、沟通和执行力，他们也可能愿意继续合作，给出暑研、RA、PhD 等 return offer。}
\SourceParagraph{另一个感受是：部分境外高校本科生能够接触到的科研资源更多。南大前 20\%的本科生能够得到的资源，可能和港三前 60\%的本科生能得到资源相似。这些资源体现在课题、指导、seminar、组会、跨组交流、访问学生制度等各个方面。境外高校的老师还可能提供更多境外 connection，利好出国申请。特别地，我相信很多境外老师会乐于接受南大的同学：他们可能曾经接收过南大学生做暑研/读 PhD，或者自己就是南大校友，更乐意接收培养南大的同学。某种意义上，这也是校友资源。}
\subsection{3. 交换与 social}
\SourceParagraph{有机会认识很棒的同学、师兄师姐和老师。office hour、coffee talk、seminar、组会、课后闲聊的机会都可以把握。进组的话可以提前体验组内生态，和师兄师姐一起吃饭，一线观察并了解一个 PhD 是怎么炼成的。Social 还是交换生活中蛮重要的一部分，所以不要错过！}
\subsection{4. 交换与其他}
\SourceParagraph{交换院校往往会提供 edu 邮箱。若你想套 HYPSM 的老师，@berkeley.edu 往往比 @nju.edu.cn 在老师的邮箱列表中更加醒目。交换 = 学 + 玩 + 工。除了上述功利因素，交换本身是一段难得的体验。交换=学+玩+工。除了上述功利因素，交换本身是一段难得的体验。交换一段时间与旅游的感觉完全不同。进入一个崭新的环境，体验当地的风土人情。相较于出国读硕士花费较少，可以提前尝试留学生活，适应全英文的生活环境。交换强迫你走出舒适区，去看到更大的世界，主动建立社交，了解各类信息和资源，使自己成为更独立的人，去应对孤独和迷茫，在陌生的地方重新整理自己的生活节奏，最终种种收获化作你的成长。私下自己觉得交换的这一学期是本科迄今为止的六个学期中收获最大的一个学期，许多观念被 refresh，许多目标被达成，很感谢当初自己选择交换一学期，某种意义上这拯救了我。}
\section{交换学习有哪些坏处}
\subsection{1. 严肃向免修不免考哈气}
\SourceParagraph{南赫有些课程只在半学期开课，如果你不幸选择这类课程免修不免考，且这类课程有考试要求（南赫哪门课没有呢，再烂的课都要考试），你将面临以下选择题，无法逃避：不回来考试的话，课程分数不会很高，甚至无法正常完成考核；回来考试则要自己承担往返路费、时间和精力。笔者交换时，学院并不允许线上考试/提交论文等方式直接替代线下考试；还有一些老师可能要求必须参与课程 pre。总之可能会有一些烦心事出现，不过 case by case，回过头来看这些也没什么}
\subsection{2. 一定！不要在！境外！搞丢你的证件！}
\SourceParagraph{护照/通行证/身份证/校园卡/学生证/宿舍门禁卡——谁会把这些都揣在自己的口袋里面然后一起弄丢呢，好难猜。你需要保管好各类证件，否则你的交换之旅将十分闹心。在内地弄丢这些证件还算可以接受，但若在境外丢失，补办和身份核验都可能变得麻烦。建议出发前给重要证件留电子扫描件或照片备份，证件和银行卡也不要全部塞在同一个包里。面对这样可怕结果的可能姑且算是交换的坏处之一吧。}
\subsection{3. “延毕风险”}
\SourceParagraph{不得不品味的一环。报名半年及以上的交换项目必然会听到这句话。延毕风险确实是需要考虑的一环，主要风险在于修不满学分、必修课安排冲突、挂科或学分认定出现问题，最终无法按培养方案正常毕业。考虑到大四之前的部分课程通常还有后续补修空间，笔者个人判断，一个对自己学业负责、提前核对培养方案的本科生，实际延毕的概率并没有想象中高。让我们期待第一届毕业生的表现。尽管如此，交换本身就是在承担一点制度和路径上的不确定性，你需要有勇气，也需要做功课。}
\subsection{4. 国际服往往是高端局（混内地人圈子 or）}
\SourceParagraph{交换意味着你会进入一个文化背景、生活习惯、边界感、饮酒习惯、dating culture、金钱观念都可能和自己熟悉环境不同的圈子。内地交换生往往比较容易沟通，你遇到的其他同学可能并非如此；与此同时，保持一点点警惕不可或缺。交友不慎可能带给你并不美妙的回忆。}
\subsection{5. 安全}
\SourceParagraph{笔者直到在 hk 的最后两周才意识到，自己交换了一个学期仍不知道香港的急救电话是什么——请不要学习笔者。出发之前，至少要知道自己前往地区的报警、急救、学校 security 等联系方式，确认自己的保险覆盖什么，也最好在当地有几个可以紧急联系的朋友/师兄师姐。平时一同出入、结伴而行、相互照应是很不错的。}
\subsection{6. 吐槽}
\SourceParagraph{吐槽一下 NJU/NHI 对交换的不支持。交换生系统上的公费交换项目往往可以免去学费，这节省了很大一笔开支，而生活费、住宿费、路费等仍需要自己支付，不同地区生活费差异很大。自费交换项目可以在交换生管理系统上申请，也可以自己到院校系统上申请（HYPSM 笔者印象中都有自费交换项目申请，具体以当年官网为准）。交换生管理系统上的自费项目主要有 UCB BISP、UPenn IGSP 等，总花费可以自己了解。很多院校会报销路费等，乃至帮助本科交换生拿到 CSC 等资助，可能覆盖交换期间相当一部分开支；但笔者交换时 NJU 并无对普通交换项目的统一资助政策，拔尖/强基等特殊班型可能会有额外支持。NHI 目前对交换本身还是挺支持的，听过数次交换宣讲，也希望未来在课程设置、免修不免考、学分衔接等方面更多地考虑交换。如果能给出更多完整、可自由支配的时间让同学们交换，而不是每学期都塞满课程，相信我们的出国去向和科研选择也会更丰富。}
\section{简单讲解一下笔者的交换经历作为参考}
\SourceParagraph{2025/8/29\hspace*{1.5em}初步确定了大三下交换的想法，在交换生平台上看到历年新二港三的交换项目都需要语言成绩，遂赶鸭子上架，临时报名了雅思考试。}
\SourceParagraph{在诸多学长学姐的推荐下，笔者很早就开始考虑交换，但是并没有提前做任何实际准备，备考语言成绩，现在来看这是败笔，语言成绩越早考出来越好。仓促决定是因为 1，大二在苏州校区感到处处受限，没有产出，十分沮丧、焦虑 2，暑假期间想好了可以申请港硕作为保底，去香港提前生活一学期，可以作为自己是否适应当地长期生活的检验 推荐任何考虑交换的同学尽快考出来语言成绩，除非考虑去 tw 这样不要求语言成绩的项目。}
\SourceParagraph{2025/9/13\hspace*{1.5em}在苏州大学考了第一次雅思（这是当时最早能够线下考试的日期，口语在 9/8 考完），考完感觉发挥并不好。}
\SourceParagraph{2025/9/24\hspace*{1.5em}报名入选了 HKUST 的交换生项目，这时候语言成绩还没出来。此时还报名了郑钢的科研项目，若两个项目都能成行，则之后一整年笔者都会在境外度过。}
\SourceParagraph{2025/9/26\hspace*{1.5em}第一次雅思出成绩，总分 6.0，两个 6.5，两个 5.5。已经预料到成绩会很差，幸而刚刚好满足了交换项目对语言成绩的要求。}
\SourceParagraph{想来也是命运使然，从 8 月末开始报名的数个交换项目中，hku 因为自己没有及时提醒院系审批而审核不通过，nus 一向竞争激烈（rank \textasciitilde{}5\%会比较有把握），意料之中落选，cuhk 要求雅思 6.5，报名后 9/5 接到了国际处的电话询问语言成绩，因为成绩还没出来，名额给了其他同学。最想去的台北最后还是未能成行。现在回头看，一切都是最好的安排。}
\SourceListItem{9 月末至 11 月上旬\hspace*{1em}笔者在学长的指点下，开始寻找 26 暑研/实习以及交换期间科研的机会，投了二十份左右的简历，大多直接石沉大海，少部分面试未通过/没有 follow up。后面通过了郑钢项目的教授面试，便不再考虑寻找新的暑研，主要开始找港校科研的机会。笔者主要找的是在做和自己 bg 比较匹配的 AI4Atmosphere 方向的老师，信息了解途径主要是微信公众号推送和学校官网的教职工介绍。目前做 AI4 大气方向的老师大致有两类，一类是大气出身，一类是 CS/AI 出身，笔者当时只套了后面一类。这一类老师数量极少，所以套磁并不顺利。}
\SourceParagraph{题外话 虽然数量不多，HKUST 还是有一些组的研究与我们传统意义上的大气科学相关，包括气候、大气化学、预报、流体等主题。其中不乏研究做得很好，并且从前还接收过南大学生做暑研或者研究生的老师，例如 Jimmy FUNG、Xiaoming SHI。此外，一些老师的研究还会涉及 earth engineering 以及 AI for SGD，详见：\url{https://wsdi.hkust.edu.hk/organization}、 \url{https://ccrs.hkust.edu.hk/}。HK 做大气的老师并不少，HKU 的席大智老师也是南大大气本科校友，学院很多老师也有 HK 的 connection。}
\SourceParagraph{2025/11/11\hspace*{1.5em}整理了之前的思路，尝试套了一位当时还在做这个方向且做得非常好的 PhD，t 哥。t 哥目前在香港读博，他很快回复了邮件，我们加了微信。}
\SourceParagraph{2025/11/20\hspace*{1.5em}t 哥介绍了做同一方向的 h 哥，在 t 哥和 h 哥的指导下，开始做一个 NWP benchmark 的工作，后面这个工作演变成了 RealBench。}
\SourceListItem{11 月下旬至 2026/1/13\hspace*{3.5em}主要在期末考试和各种 pre 的枪林弹雨中辗转求生，期间保持着一周一次的线上 meeting 指导，并一步一步完成选宿舍、选课等任务。期末结束后开始准备通行证等证件，因为时间紧迫，联系导员帮忙开了一份加急证明，顺利在出发前取得了证件。感谢行政老师的帮助。}
\SourceParagraph{2026/1/26\hspace*{2.0em}苏州 -> 上海 -> 深圳 -> 香港。t 哥春季学期外出实习，让我坐他空出来的工位，就在 h 哥旁边。借着地理位置的便利，请教了 h 哥很多问题，手头的项目也逐渐走上了正轨，一点一点向前推进。在 t 哥和 h 哥的帮助和鼓励下，我开始旁听一周一次的大组会，为了将项目发表出去，以及要到大老板的推荐信而努力。}
\SourceParagraph{HKUST 依山面海，坐拥一流风光，无愧于“亚洲最美大学”的美誉。我平时很少上课，中午吃完饭后回宿舍小憩，下午和晚上都在工位干活。每天我还会和 h 哥及其他 PhD 师兄一起吃晚饭，一起回宿舍，听师兄师姐闲聊，学习 PhD 生存指南。从如何使用 AI、如何高效产出，到如何松弛、如何自洽，每一天都有新的感悟，每一天都有十足的收获。尽管如此，很多深夜里，我开始更深地焦虑迷茫。随着日子流逝，紧迫感渐渐涌上心头。情况与我最开始想的确乎不同，读博士与否，科研方向也罢，我总是在被动地前进。南北多歧路，无为泪沾巾。我开始发现，自己的焦虑并不是来自那些自以为是的信念，而是我的“无知”。幸运的是，我多么擅长、多么习惯找到、认定自己的“错误”，然后鞭笞它，从而领会到，犯错也没那么可怕。做一个不可救药的乐观主义者，相信未来会更好。后面的路，就是每个人自己的体验了。之后剩下的部分，我会在飞跃手册中补充。}
\section{By the end}
\subsection{你应该选择什么时候出去交换？}
\SourceParagraph{考虑长期项目的话，从课业压力和本科规划来说，大二下 / 大三上 / 大三下通常都是不错的选项。原因很现实：这几个学期如果留在学校，你可能要背负很大的学业压力，无暇去做那些真正有意义、能够帮你抵达自己上限的事情，例如科研；又或者，你正好处在心灵迷茫、疲惫的时候。交换不一定能解决你的问题，但它会强迫你换一个环境重新生活。对我而言，交换是本科阶段绝无仅有的美好体验，也可能是人生中少数可以以学生身份在境外长住、低成本体验另一种社会和学校的机会。请认真考虑、尽量争取。}
\subsection{你应该如何选择自己的交换项目？}
\SourceParagraph{很多参考，包括 bar、地区、学校档次、课程自由度、科研机会、语言环境、生活成本、是否想留在当地读研/读博等等。如果你目标很明确，比如为了科研，那就优先看能不能进组、有没有匹配的老师；如果你只是想看看世界，那地区和生活体验可以比 ranking 更重要。不需要把一切想得太清楚。就像笔者一波三折的交换经历中，最后的去向总是多种偶然共同作用的结果。大家都会去到自己该去的地方的。}
\section{Summer School / 暑研}
\SourceParagraph{短期项目里，summer school 更偏课程与旅游体验；暑期科研则对于出国申请尤其重要，其价值可能在于推荐信 \&\& 海外经历 \&\& return offer \&\& paper，也是“connection is all you need”的世界里建立联系最好的途径之一。暑研有官方渠道和自己申请的区分。官方渠道可以关注 Mitacs、UC Davis 等项目；自己申请则既有境外院校统一渠道（自费或带奖），例如 Caltech SURF、西湖暑研、台大暑研等，也可以自己套老师。交换系统上的郑钢奖学金项目比较特殊，是非常好的项目，推荐大家了解。}
\section{写在最后}
\SourceParagraph{过去的永远过去了，这是注定的，所以不要搞砸现在，好好做手头的事。及时当勉励，岁月不待人。交换给予本科生一次离开固有路径、重新理解自己和世界的机会。期望我们成为更广阔的自己，打开世界的轮廓。期待看到更多的 NHI 同学报名交换。}

